\section{Background and Preliminaries}\label{sec:background}
% TODO: Define notation, basic concepts, and any baseline models needed for understanding.

Current neuronal network architectures predominantly rely on matrix multiplications to process and transform data. This section provides an overview of the fundamental concepts and notations that underpin traditional neural networks, as well as the preliminaries necessary for understanding the proposed bit-level operations.

\subsection{Notation}
We denote vectors in bold lowercase (e.g., \(\mathbf{x}\)) and matrices in bold uppercase (e.g., \(\mathbf{W}\)). The \(i\)-th element of a vector \(\mathbf{x}\) is denoted as \(x_i\), and the element in the \(i\)-th row and \(j\)-th column of a matrix \(\mathbf{W}\) is denoted as \(W_{ij}\).

\subsection{Basic Concepts}
Neural networks consist of layers of interconnected nodes (neurons), where each connection is associated with a weight. The output of each neuron is computed as a weighted sum of its inputs, followed by a non-linear activation function. Formally, for a single neuron, the output \(y\) can be expressed as:
\[
y = f\left(\sum_{i} W_{i} x_{i} + b\right)
\]
where \(f\) is the activation function and \(b\) is the bias term.

\subsection{Baseline Models}
The most common baseline model for comparison in this paper is the fully connected feedforward neural network (FNN), which consists of multiple layers of neurons with dense connections. Other relevant models include convolutional neural networks (CNNs) for image processing tasks and recurrent neural networks (RNNs) for sequential data.
